\chapter*{Liste des acronymes et abréviations}
\addcontentsline{toc}{chapter}{Liste des acronymes et abréviations}

\begin{longtable}{p{3cm}p{11cm}}
\toprule
\textbf{Sigle} & \textbf{Signification} \\
\midrule
API & Application Programming Interface. \\
BGL & BlueGene/L, jeu de données de journaux système issu d'un supercalculateur. \\
CEF & Common Event Format. \\
CPU & Central Processing Unit. \\
CSV & Comma-Separated Values. \\
F1-score & Mesure harmonique combinant précision et rappel. \\
FastAPI & Framework Python utilisé pour exposer l'API locale du prototype. \\
FN & False Negative, anomalie réelle non détectée par le modèle. \\
FP & False Positive, événement normal signalé à tort comme anomalie. \\
HDFS & Hadoop Distributed File System. \\
HIDS & Host-based Intrusion Detection System. \\
IA & Intelligence artificielle. \\
IDS & Intrusion Detection System. \\
JSONL & JSON Lines, format de stockage ligne par ligne. \\
LEEF & Log Event Extended Format. \\
\logminer & Prototype réalisé dans ce travail pour l'analyse de journaux hétérogènes. \\
MCC & Matthews Correlation Coefficient. \\
MITRE ATT\&CK & Base de connaissance des tactiques et techniques adverses. \\
MQTT & Message Queuing Telemetry Transport. \\
PR-AUC & Aire sous la courbe précision--rappel. \\
RAM & Random Access Memory. \\
Redis Streams & Mécanisme de flux Redis utilisé comme extension événementielle locale. \\
RF & Random Forest. \\
SIEM & Security Information and Event Management. \\
SOC & Security Operations Center. \\
TN & True Negative, événement normal correctement reconnu. \\
TP & True Positive, anomalie correctement détectée. \\
Wazuh & Plateforme open source de supervision sécurité et HIDS/SIEM. \\
\bottomrule
\end{longtable}

